\documentclass[quick,pro]{debate}
\motion{短视频提升了 / 降低了人的认知能力}
\stance{正方}{短视频提升了人的认知能力}
\begin{document}
\debatetitle
\begin{thesisbox}[一句话立论]短视频第一次把知识送到了十亿普通人手里，它带来的理解增量，远大于它拿走的那点专注。\end{thesisbox}
\section{定义与判准（标准表述，全队同一口径）}
\begin{fields}
\field{短视频}{以手机为主要终端、单条时长通常在数分钟以内、由算法以信息流方式连续分发的视频内容，涵盖娱乐、生活、知识科普。
}
\field{认知能力}{知觉、注意、记忆、理解、推理、判断与语言这一整套心理加工能力（美国心理学会《心理学词典》）。\textbf{不只是专注力。}
}
\field{判准}{在这一整套维度上，比较普通人在短视频普及前后的整体状态，看\textbf{净方向}；并且要看多数人的典型使用，要分离短视频本身的贡献。\textbf{我方要证明}增益压过损耗；\textbf{对方要证明}损耗压过增益。\textbf{只证明“有好处/有坏处”都不算完成任务。}
}
\field{对方提出偏向定义或判准时的 30 秒口径}{对方若把短视频定义成“算法投喂的娱乐内容”——“你把结论写进了定义，知识科普是平台的正式内容分类，不能靠定义开除。”对方若主张“脱离设备时的能力才算”——“按这个标准，书本、笔记、计算器都在降低人的认知能力。这不是判准，是预设结论。”
}
\end{fields}
\section{我方论点}
\begin{longtable}{@{}L{\dimexpr0.047\linewidth-1.50\tabcolsep}L{\dimexpr0.187\linewidth-1.50\tabcolsep}L{\dimexpr0.466\linewidth-1.50\tabcolsep}L{\dimexpr0.301\linewidth-1.50\tabcolsep}@{}}
\toprule \thead{标签} & \thead{一句话} & \thead{核心数据 / 学理 / 案例} & \thead{出处}\\\midrule\endhead
\textbf{知识普惠} & 短视频把知识的门槛降到零，送到了十亿量级普通人面前，群体的知识与理解被整体抬高 & 数据：截至 2024 年 12 月，短视频用户 10.40 亿人，占网民整体 93.8\%，其中农村网民 3.13 亿人。学理：晶体智力靠后天接触积累、随经验增长。案例：2024 年 5 月抖音重点支持一百位院士入驻做科普 & CNNIC 第 55 次报告（已核实）；Cattell 1963（有把握）；媒体报道（有把握）\\
\textbf{形式适配} & 短 + 看得见听得见 + 能自控节奏，本身就是有效的认知形式 & 数据：Guo 等对 edX 四门课程 690 万次观看会话的分析，视频越短参与度越高。学理：Mayer 分段原则与双通道假设。案例：可汗学院以短视频起家 & Guo, Kim \& Rubin 2014, ACM L@S（已核实）；Mayer《Multimedia Learning》（有把握）\\
\bottomrule\end{longtable}
\section{对方论点 → 一句话反驳}
\begin{longtable}{@{}L{\dimexpr0.163\linewidth-1.33\tabcolsep}L{\dimexpr0.352\linewidth-1.33\tabcolsep}L{\dimexpr0.485\linewidth-1.33\tabcolsep}@{}}
\toprule \thead{对方会说} & \thead{我方一句话回} & \thead{对方追问时}\\\midrule\endhead
注意力损耗（脑电研究） & 48 人的横断相关，作者自己写着不能谈因果 & “你承认注意力受影响吗？”→“承认现象值得关心；不承认因果，也不承认它能翻掉总账。”\\
时段碎片化 & 你没证明这些时间原本用在认知价值更高的事上 & “我说的是切碎不是换走”→“那请证明：一，人原本有整块时间；二，切碎它的是短视频而不是所有推送。”\\
懂了的错觉 & 纪录片、科普书都有这个问题，你要证明短视频特别严重 & “评论区也是错的”→“多个声音互相冲突，正是判断力的练习条件。”\\
算法替你筛选 / 信息茧房 & 多数人过去的信息来源是电视和熟人，主动性同样是零；信息流的问题是杂不是窄 & “那不就交给算法了？”→“你在把‘谁替你选’当成‘你会不会想’。”\\
\bottomrule\end{longtable}
\section{必问与必答}
\begin{points}{必问 （对方不答就点名，自由辩前两分钟内问完）}
\item \textbf{你方证明的是有损伤，还是净降低？}（全场核心）
\item 48 人的横断相关，能推到 10.40 亿人吗？
\item 短视频在“让人分心”里占多少份额？你方分离过吗？
\item \textbf{短视频之前，那 3.13 亿农村网民从哪里获得这些知识？}（一辩稿已抛出，全场追问）
\end{points}
\begin{points}{必答 （15 秒内正面回）}
\item 承认注意力受影响吗 → “承认现象值得关心，不承认因果被证明，也不承认它足以翻转总账。”
\item 看到等于学到吗 → “不等于。我方主张的是接触量的量级跃迁作用于晶体智力这一维度。”
\item 数据是平台发的吗 → “其中一份是中国科普研究所与抖音联合发布，我方一辩就交代了；覆盖面数据来自 CNNIC。”
\item 短视频里大部分是娱乐吧 → “双方今天都没有占比数字，所以谁都不能用它论证。”
\end{points}
\section{自由辩战场概要}
\begin{longtable}{@{}L{\dimexpr0.136\linewidth-1.50\tabcolsep}L{\dimexpr0.299\linewidth-1.50\tabcolsep}L{\dimexpr0.252\linewidth-1.50\tabcolsep}L{\dimexpr0.313\linewidth-1.50\tabcolsep}@{}}
\toprule \thead{战场} & \thead{开场问题} & \thead{目标} & \thead{收束语}\\\midrule\endhead
一 · 净方向 & 你方证明的是有损伤，还是净降低？ & 让评委记住对方只证明了存在 & “这道题问的是总账。总账，对方今天没有算过。”\\
二 · 归因与证据强度 & 让人分心的东西那么多，短视频占多少份额？ & 让评委觉得对方证据撑不住结论 & “双方都只有相关时，评委该看量级和覆盖面。”\\
三 · 谁被接住了 & 短视频之前，3.13 亿农村网民从哪里获得知识？ & 把直觉从“我浪费了时间”转到“有人本来什么都没有” & “这道题不能只替坐在图书馆里的人算。”\\
\bottomrule\end{longtable}
时间分配：前 90 秒战场一（抛完四个必问）；中 90 秒战场二；后 60 秒战场三收束，不开新战场；留 20–30 秒备用。站位：正一守定义与判准，正二接证据强度，正三接盘问成果，正四收束与价值。

\section{如果……就……}
\begin{contingency}
\ifthen{对方定义不同，或拿出没预料到的数据}{定义：“你把结论写进了定义”，随即回到中立定义原话，不纠缠。新数据：先问机构与年份，再问“它测的是‘看到’还是‘能力’？”}
\ifthen{论点二被打穿，或对方回避必问}{收缩到论点一：“就算形式这一层打平，覆盖面的量级差还在。”对方回避就点名、重复、不换题。}
\ifthen{纪律}{不说“刷短视频等于学习”；不说“93.8\% 的中国人”（分母是网民）；不说未核实的效应量数字。}
\end{contingency}
\end{document}
